\rtmtight
\begin{tblr}{width=\textwidth,colspec={|Q[c,m,2.1cm]|X[l,m]|},hlines,vlines,hspan=minimal,rowsep=1.5pt,colsep=3pt,row{1}={bg=headgray,font=\bfseries}}
\SetCell{c,m}\hfil\logo\hfil & \textbf{POLITEKNIK NEGERI MALANG}\par\textbf{JURUSAN TEKNOLOGI INFORMASI}\par\textbf{PROGRAM STUDI: D4 TEKNIK INFORMATIKA} \\
\SetCell[c=2]{l} \textbf{RENCANA TUGAS MAHASISWA (RTM) DAN RUBRIK PENILAIAN} & \\
Nama Mata Kuliah & Desain dan Pemrograman Web \\
Kode Mata Kuliah & RTI253004 \quad \textbf{sks / Jam perminggu:} 3 SKS/6 jam \quad \textbf{Semester:} 3 \\
Dosen Pengampu & \begin{enumerate}\item Priska Choirina, M.Tr.T.\item Muhammad Unggul Pamenang, S.St., M.T.\item Dimas Wahyu Wibowo, S.T., M.T.\item Farid Angga Pribadi, S.Kom., M.Kom.\end{enumerate} \\
\end{tblr}
\assessmentblock{Tugas HTML, CSS, JavaScript, dan PHP}{Hands-on Practice}{SCPMK0704-02201}{Mahasiswa menyelesaikan tugas html, css, javascript, dan php sesuai Sub-CPMK dan konteks mata kuliah.}{Menganalisis kebutuhan, mengerjakan artefak praktik/proyek, menguji hasil, mendokumentasikan evidence, dan mempresentasikan keputusan bila diminta.}{Laporan, artefak digital, repository/prototype, evidence pengujian, dan/atau bahan presentasi.}{Struktur halaman dan semantic HTML & 9.6 & Halaman memakai struktur semantic, form, navigasi, dan konten yang sesuai kebutuhan fitur. & Struktur halaman cukup benar, tetapi semantic element atau kelengkapan form masih terbatas. & HTML tidak terstruktur atau tidak mendukung kebutuhan fitur. \\ \\
Kualitas styling dan interaksi JavaScript & 8.4 & CSS responsif dan konsisten, JavaScript menangani interaksi/validasi sisi klien dengan jelas. & Tampilan dan interaksi berjalan, tetapi responsivitas atau validasi masih kurang. & Styling berantakan atau JavaScript tidak berjalan sesuai kebutuhan. \\ \\
Integrasi PHP dasar dan evidence & 6 & PHP memproses input/output dengan benar, validasi dasar tersedia, dan hasil diuji serta terdokumentasi. & PHP berjalan, tetapi validasi, error handling, atau evidence pengujian belum lengkap. & PHP tidak memproses data dengan benar atau tidak ada bukti pengujian. \\}

\assessmentblock{UTS Modul Web Dasar}{UTS praktik}{SCPMK0704-02201, SCPMK0209-02202}{Mahasiswa menyelesaikan uts modul web dasar sesuai Sub-CPMK dan konteks mata kuliah.}{Menganalisis kebutuhan, mengerjakan artefak praktik/proyek, menguji hasil, mendokumentasikan evidence, dan mempresentasikan keputusan bila diminta.}{Laporan, artefak digital, repository/prototype, evidence pengujian, dan/atau bahan presentasi.}{Kelengkapan modul web dasar & 8 & Modul mencakup halaman, form, validasi, proses PHP, dan alur navigasi sesuai spesifikasi. & Modul utama tersedia, tetapi ada bagian form, validasi, atau navigasi yang belum lengkap. & Modul tidak memenuhi fitur dasar yang diminta. \\ \\
Kerapian kode dan pemisahan concern & 7 & HTML, CSS, JavaScript, dan PHP tertata, mudah dibaca, dan tidak bercampur secara berlebihan. & Kode cukup berjalan, tetapi struktur file atau pemisahan concern belum rapi. & Kode sulit dipelihara atau banyak duplikasi tanpa struktur. \\ \\
Pengujian dan penjelasan hasil & 5 & Kasus uji input normal/invalid dilakukan, output diverifikasi, dan hasil dapat dijelaskan saat praktik. & Pengujian dilakukan terbatas, atau penjelasan hasil belum lengkap. & Tidak ada pengujian bermakna atau tidak mampu menjelaskan hasil modul. \\}

\assessmentblock{Tugas PHP, MySQL, Session, dan Hosting}{Hands-on Practice}{SCPMK0209-02202, SCPMK0603-02203}{Mahasiswa menyelesaikan tugas php, mysql, session, dan hosting sesuai Sub-CPMK dan konteks mata kuliah.}{Menganalisis kebutuhan, mengerjakan artefak praktik/proyek, menguji hasil, mendokumentasikan evidence, dan mempresentasikan keputusan bila diminta.}{Laporan, artefak digital, repository/prototype, evidence pengujian, dan/atau bahan presentasi.}{Desain database dan operasi CRUD & 8 & Skema MySQL, relasi sederhana, query, dan operasi CRUD sesuai kebutuhan data aplikasi. & CRUD berjalan, tetapi skema, validasi query, atau relasi masih kurang konsisten. & CRUD gagal atau desain database tidak sesuai kebutuhan. \\ \\
Session, autentikasi, dan alur aplikasi & 7 & Session, login/logout, proteksi halaman, dan alur pengguna diterapkan benar. & Session berjalan, tetapi proteksi halaman atau alur login masih memiliki celah. & Session/autentikasi tidak bekerja atau halaman sensitif tidak terlindungi. \\ \\
Deployment hosting dan evidence uji & 5 & Aplikasi berhasil dihosting, konfigurasi koneksi benar, dan evidence uji fitur online terdokumentasi. & Deployment berhasil sebagian, tetapi konfigurasi atau evidence pengujian belum lengkap. & Aplikasi tidak dapat diakses di hosting atau fitur utama gagal. \\}

\assessmentblock{PBL Aplikasi Web Dinamis}{Project Based Learning}{SCPMK0603-02203, SCPMK0603-02204}{Mahasiswa menyelesaikan pbl aplikasi web dinamis sesuai Sub-CPMK dan konteks mata kuliah.}{Menganalisis kebutuhan, mengerjakan artefak praktik/proyek, menguji hasil, mendokumentasikan evidence, dan mempresentasikan keputusan bila diminta.}{Laporan, artefak digital, repository/prototype, evidence pengujian, dan/atau bahan presentasi.}{Kesesuaian fitur dengan kebutuhan pengguna & 4.4 & Fitur web dinamis dirancang dari kebutuhan kasus, prioritas jelas, dan alur pengguna lengkap. & Fitur utama tersedia, tetapi prioritas atau hubungan dengan kebutuhan pengguna belum kuat. & Fitur tidak menjawab kebutuhan kasus atau alur pengguna tidak jelas. \\ \\
Kualitas implementasi frontend-backend-database & 3.9 & Frontend, PHP backend, database, validasi, session, dan error handling terintegrasi stabil. & Integrasi berjalan, tetapi validasi, error handling, atau struktur kode masih terbatas. & Komponen utama tidak terintegrasi atau aplikasi sering gagal. \\ \\
Pengujian, deployment, dan presentasi proyek & 2.7 & Fitur diuji, aplikasi dideploy, evidence lengkap, dan presentasi menjelaskan keputusan teknis. & Pengujian/deployment tersedia, tetapi evidence atau presentasi keputusan masih kurang. & Tidak ada pengujian/deployment memadai atau proyek tidak dapat didemonstrasikan. \\}

\assessmentblock{UAS Demonstrasi Aplikasi Web}{UAS praktik dan defense}{SCPMK0603-02203, SCPMK0603-02204}{Mahasiswa menyelesaikan uas demonstrasi aplikasi web sesuai Sub-CPMK dan konteks mata kuliah.}{Menganalisis kebutuhan, mengerjakan artefak praktik/proyek, menguji hasil, mendokumentasikan evidence, dan mempresentasikan keputusan bila diminta.}{Laporan, artefak digital, repository/prototype, evidence pengujian, dan/atau bahan presentasi.}{Kelengkapan demonstrasi fitur & 10 & Demonstrasi mencakup alur utama, validasi input, CRUD, session, dan kondisi error. & Demonstrasi menunjukkan fitur utama, tetapi beberapa validasi atau kondisi error belum ditunjukkan. & Demonstrasi tidak membuktikan fitur utama berjalan. \\ \\
Kesiapan aplikasi dan kualitas teknis & 8.8 & Aplikasi stabil, struktur kode rapi, database konsisten, dan deployment siap diakses. & Aplikasi cukup stabil, tetapi struktur kode, database, atau deployment masih memiliki kekurangan. & Aplikasi tidak stabil atau gagal saat demonstrasi. \\ \\
Defense dan rekomendasi perbaikan & 6.2 & Mahasiswa menjelaskan arsitektur, keputusan teknis, limitasi, dan rencana perbaikan berbasis evidence. & Defense cukup jelas, tetapi limitasi atau rekomendasi belum mendalam. & Tidak mampu menjelaskan aplikasi atau rekomendasi tidak relevan. \\}

\begin{tblr}{width=\textwidth,colspec={|Q[l,m,3.2cm]|X[l,m]|},hlines,vlines,hspan=minimal,row{1-3}={bg=headgray},rowsep=1.5pt,colsep=3pt}
Jadwal Pelaksanaan: & Minggu 1--17 sesuai RPS. Setiap evaluasi memiliki RTM dan rubrik multi-kriteria. \\
Lain-Lain Yang Diperlukan: & LMS, repositori/prototype/proyek, perangkat lunak sesuai mata kuliah, dan perangkat presentasi. \\
Pustaka: & Mengacu pustaka utama dan pendukung pada bagian RPS. \\
\end{tblr}
